\section{Related Work}
\label{sec:related}

Three converging research threads form the prior art for \sysname.
We summarise each, then state the gap.

\smallskip
\noindent\textbf{Consent management on distributed ledgers.}
MedRec~\cite{azaria2016medrec} pioneered the use of a blockchain
substrate to record patient consent for medical-data access and
permissioning. The system established the durability and auditability
properties we inherit: every consent state transition is an immutable
log entry, and the consent owner retains exclusive control over
modifications. MedRec, however, treats the consent record as the
\emph{endpoint} of the protocol --- there is no downstream ML training
whose execution is gated by the consent state, and no mechanism for
demonstrating that withdrawn consent translated into a model update.
Subsequent semantic-web work such as GConsent and the W3C DPV
\cite{pandl2021gconsent} formalises consent vocabularies but
operates in an off-chain registry without enforcement.
\emph{Gap}: consent is recorded but does not influence training.

\smallskip
\noindent\textbf{Machine unlearning.}
The SISA training scheme~\cite{bourtoule2021sisa} pre-shards training
data and caches per-shard intermediate models so that data removal
requires retraining only the affected shard. FedEraser
\cite{liu2021federaser} extends SISA's logic to the federated setting
by replaying aggregation history with a calibration step,
significantly cheaper than full retraining. Gradient-ascent unlearning
methods~\cite{nguyen2022survey} take a single trained model and
perform a small number of negative-learning-rate steps on the data to
be removed. All three families have been validated on standalone
benchmarks but \emph{none specifies how unlearning is triggered, by
whom, or how a deployer demonstrates it actually happened}. The
unlearning literature treats the trigger and verification steps as
out-of-scope.
\emph{Gap}: unlearning is implementable but not verifiable
end-to-end.

\smallskip
\noindent\textbf{Blockchain-based federated learning.}
BlockFL~\cite{kim2020blockchained} records every device-side model
update on-chain to defend against malicious aggregators. BFLC
\cite{li2020bflc} replaces the FL aggregator with a committee-consensus
mechanism. These systems address aggregator-side trust but do not
encode consent at all --- a participant who joins the federation
cannot subsequently withdraw without trusting the orchestrator to
respect the request. Membership inference attacks
\cite{shokri2017mia,carlini2022lira,choquettechoo2021labelonly} provide
the canonical operational test for whether a model still memorises a
specific data subject's training contribution, but no published
blockchain-FL system uses MIA scores as the on-chain acceptance
criterion for a state transition.
\emph{Gap}: model updates are auditable but consent semantics and
forgetting verification are absent.

\smallskip
\noindent\textbf{Positioning.}
\sysname is the first system, to our knowledge, that binds all three
threads. We adopt MedRec's on-chain consent registry pattern (vertex
A), extend it with the \clsm safety and liveness properties; we accept
any of the three established unlearning paradigms as a black-box
(vertex B); and we make the verification step (vertex C) the chain's
gating condition for a consent state transition. The pattern that
emerges --- consent change $\rightarrow$ unlearning execution
$\rightarrow$ three-MIA verification $\rightarrow$ on-chain attestation
$\rightarrow$ consent state finalisation --- is the missing edge.

The choice of \emph{which} blockchain matters at IoT scale.
Permissionless gas-based chains (Ethereum, Polygon) impose recurring
per-operation costs that scale linearly with device count and consent
operation rate. Section~\ref{sec:evaluation} quantifies this against
PureChain, a public, gas-free, EVM-compatible chain with PoA$^2$
consensus and Merkle-proof-backed state. PureChain is not the
contribution of this paper; rather, our contribution is showing that
the protocol is operationally identical on both chain types, while
the economic model is fundamentally different.
